\documentclass[11pt,a4paper]{article}

\usepackage{fontspec}
\setmainfont{DejaVu Serif}[
  BoldFont={DejaVu Serif Bold},
  ItalicFont={DejaVu Serif},
  ItalicFeatures={FakeSlant=0.18},
  BoldItalicFont={DejaVu Serif Bold},
  BoldItalicFeatures={FakeSlant=0.18}
]
\setsansfont{DejaVu Sans}
\usepackage{polyglossia}
\setmainlanguage{english}
\usepackage[a4paper,margin=21mm]{geometry}
\usepackage{microtype}
\usepackage{setspace}
\setstretch{1.04}
\usepackage{amsmath,amssymb,amsthm,mathtools,bm,mathrsfs}
\usepackage{booktabs,longtable,array,tabularx}
\usepackage{enumitem}
\usepackage{tocloft}
\setlength{\cftsecnumwidth}{2.5em}
\setlength{\cftsubsecnumwidth}{3.4em}
\usepackage{xcolor}
\usepackage{tikz}
\usetikzlibrary{arrows.meta,positioning,calc,fit}
\usepackage[unicode,hidelinks]{hyperref}
\usepackage[nameinlink,noabbrev]{cleveref}
\emergencystretch=3em
\allowdisplaybreaks

\hypersetup{
  pdftitle={Tensorial Observation Geometry in General Relativity},
  pdfauthor={Stassis Research Program},
  pdfsubject={GO-QS reference document for coordinate covariance, frames, pullbacks, and tensor descent},
  pdfkeywords={general relativity, tetrad, coframe, observation map, pullback, basic tensor, GO-QS}
}

\definecolor{obsblue}{RGB}{37,83,138}
\definecolor{obsred}{RGB}{150,45,45}
\definecolor{obsgreen}{RGB}{30,115,78}
\definecolor{lightblue}{RGB}{239,246,253}
\definecolor{lightred}{RGB}{253,242,242}

\theoremstyle{definition}
\newtheorem{definition}{Definition}[section]
\newtheorem{model}[definition]{Model}
\theoremstyle{plain}
\newtheorem{theorem}[definition]{Theorem}
\newtheorem{proposition}[definition]{Proposition}
\newtheorem{lemma}[definition]{Lemma}
\newtheorem{corollary}[definition]{Corollary}
\theoremstyle{remark}
\newtheorem{remark}[definition]{Remark}
\newtheorem{warning}[definition]{Boundary}

\newcommand{\R}{\mathbb{R}}
\newcommand{\B}{\mathcal{B}}
\newcommand{\F}{\mathcal{F}}
\newcommand{\Qreg}{\mathcal{Q}}
\newcommand{\Dprof}{\mathscr{D}}
\newcommand{\dd}{\mathop{}\!d}
\newcommand{\norm}[1]{\left\lVert #1\right\rVert}
\newcommand{\abs}[1]{\left\lvert #1\right\rvert}
\newcommand{\set}[1]{\left\{#1\right\}}
\newcommand{\given}{\,\middle|\,}
\newcommand{\rad}{\operatorname{rad}}
\newcommand{\im}{\operatorname{im}}
\newcommand{\rank}{\operatorname{rank}}
\newcommand{\id}{\operatorname{id}}
\newcommand{\SO}{\operatorname{SO}}
\newcommand{\Lie}{\mathcal{L}}
\newcommand{\status}[1]{\textsf{\footnotesize #1}}

\newcolumntype{L}[1]{>{\raggedright\arraybackslash}p{#1}}
\newcolumntype{Y}{>{\raggedright\arraybackslash}X}

\title{\bfseries Tensorial Observation Geometry in General Relativity\\
\large Coordinate covariance, frame/coframe duality, observation maps, and tensor descent}
\author{Stassis Research Program}
\date{Strict reference revision v0.3 (P0) --- 2026}

\begin{document}
\maketitle

\begin{center}
\fcolorbox{obsblue}{lightblue}{%
\parbox{0.91\linewidth}{\small
\textbf{Status.} First tensor/GR reference document migrated to GO-QS v0.1.
It corrects the tetrad reconstruction formula, declares the unit convention,
states exact hypotheses for tensor descent, and separates coordinate,
gauge, apparatus, and observation transformations. It does not modify
Einstein's equations or claim a new theory of gravity.}}
\end{center}

\begin{abstract}
This note supplies a type-safe tensorial layer for Geometry of Observation.
Four operations are separated: passive coordinate re-description, active
diffeomorphism, local Lorentz frame gauge, and a possibly non-injective
observation channel. Frame coefficients $e_a{}^\mu$ and coframe coefficients
$e^a{}_\mu$ are distinguished; the covariant metric is reconstructed from
the coframe, while the inverse metric is reconstructed from the frame.
The default convention is signature $(-,+,+,+)$ and geometrized units
$c=G=1$, with the SI conversion stated explicitly. For a smooth observation
map, the radical of a pulled-back pseudo-Riemannian form is computed exactly,
showing why nontrivial kernel is sufficient but not necessary for degeneracy
in Lorentzian signature. Smooth scalar descent and descent of covariant
tensors along a surjective submersion are stated with their required
fiber hypotheses. A local Lorentz gauge change is distinguished from a
physical reorientation of an apparatus. A GO-QS quantity and protocol ledger
closes the document.
\end{abstract}

\tableofcontents

\section{Scope, status, and conventions}

Unless stated otherwise, manifolds are smooth, Hausdorff, and second
countable. A spacetime $(M,g)$ is a four-dimensional, time-oriented
Lorentzian manifold with signature
\[
\eta_{ab}=\operatorname{diag}(-1,1,1,1).
\]
The curvature convention is
\[
R(X,Y)Z
=\nabla_X\nabla_YZ-\nabla_Y\nabla_XZ-\nabla_{[X,Y]}Z.
\]
All relativistic formulae use the GO-QS unit context
\texttt{GR\_geometrized}: $c=G=1$. Hence a proper-time parametrized
four-velocity satisfies $g(u,u)=-1$. If proper time is measured in SI
seconds and $u_{\rm SI}=\dd\gamma/\dd\tau$, then
\[
g(u_{\rm SI},u_{\rm SI})=-c^2.
\]

\begin{warning}[Claim boundary]
The tensor calculus, frame bundle, Fermi--Walker transport, and submersion
results used below are standard. The contribution here is a corrected,
typed interface to observation geometry. No new field equation, force,
or gravitational degree of freedom is introduced.
\end{warning}

\section{Four transformations that must not be conflated}

\subsection{Passive coordinate change}

Let $x:U\to\R^4$ and $x':U'\to\R^4$ be charts. On the overlap,
\[
x'\circ x^{-1}:x(U\cap U')\to x'(U\cap U')
\]
is a diffeomorphism between coordinate images. It changes components, not
the geometric point or tensor. With
\[
J^\alpha{}_\mu=\frac{\partial x'^\alpha}{\partial x^\mu},
\qquad
(J^{-1})^\mu{}_\alpha
=\frac{\partial x^\mu}{\partial x'^\alpha},
\]
the basis and components transform contragrediently.

\subsection{Active diffeomorphism}

An active diffeomorphism $\varphi:M\to M$ maps fields by pullback or
pushforward. Whether it represents a gauge redundancy, a physical symmetry,
or a change of boundary data depends on the theory and boundary conditions.
It is not identified with a passive chart change merely because both use
Jacobians.

\subsection{Local Lorentz frame gauge}

For an orthonormal frame $e_a$, a local Lorentz transformation is
\[
e'_a=\Lambda_a{}^b e_b,
\qquad
\Lambda(x)\in\SO^+(1,3).
\]
It changes the frame representation. If all frame components of physical
fields are transformed consistently, gauge-invariant readouts do not change.

\subsection{Observation and reduction}

GO-QS assigns distinct signatures:
\[
T_g:E\to E,\qquad
R_\lambda:E\to E_\lambda,\qquad
O_\lambda:E_\lambda\to Y_\lambda^0,\qquad
K_\lambda(\dd y\mid y^0).
\]
$T_g$ is invertible; $R_\lambda$ and $O_\lambda$ need not be. Information
loss is therefore a property of a declared channel, not of covariance
itself.

\begin{figure}[htbp]
\centering
\begin{tikzpicture}[
  node distance=8mm and 12mm,
  box/.style={draw,rounded corners,minimum height=9mm,align=center,fill=lightblue},
  arr/.style={-{Latex[length=2.2mm]},thick}
]
\node[box] (geom) {geometric fields\\on $M$};
\node[box,right=of geom] (repr) {chart/frame\\components};
\node[box,below=of geom] (red) {reduced state\\$E_\lambda$};
\node[box,right=of red] (data) {sensor data\\$Y_\lambda$};
\draw[arr,<->] (geom) -- node[above] {representation} (repr);
\draw[arr] (geom) -- node[left] {$R_\lambda$} (red);
\draw[arr] (red) -- node[above] {$O_\lambda,K_\lambda$} (data);
\draw[arr,obsred,bend left=35] (geom) to node[left] {$\varphi$} (geom);
\end{tikzpicture}
\caption{Representation changes preserve the geometric object; an observation
pipeline can quotient or randomize it.}
\end{figure}

\section{Tensors and coordinate covariance}

\begin{definition}[Tensor]
At $p\in M$, a tensor of type $(r,s)$ is a multilinear map
\[
T:(T_p^*M)^r\times(T_pM)^s\to\R.
\]
\end{definition}

For $V=V^\mu\partial_\mu$ and
$\omega=\omega_\mu\dd x^\mu$,
\[
V'^\alpha
=\frac{\partial x'^\alpha}{\partial x^\mu}V^\mu,
\qquad
\omega'_\alpha
=\frac{\partial x^\mu}{\partial x'^\alpha}\omega_\mu.
\]
Consequently, the contraction $\omega_\mu V^\mu$ is coordinate
independent. General tensor components receive one Jacobian per upper
index and one inverse Jacobian per lower index.

\begin{remark}[Coordinate passport]
A component array is not a complete GO-QS quantity. It must carry the
chart, basis, coordinate dimensions, index variance, and frame/gauge law.
If coordinates have different physical dimensions, component dimensions
are index dependent.
\end{remark}

\section{Metric, connection, curvature, and basepoint}

The metric is
\[
g=g_{\mu\nu}\,\dd x^\mu\otimes\dd x^\nu,
\qquad
\dd s^2=g_{\mu\nu}\dd x^\mu\dd x^\nu.
\]
Under a coordinate change,
\[
g'_{\alpha\beta}
=\frac{\partial x^\mu}{\partial x'^\alpha}
\frac{\partial x^\nu}{\partial x'^\beta}g_{\mu\nu},
\]
and the quadratic form $\dd s^2$ is unchanged. In Lorentzian signature,
this quadratic form is not a metric distance function.

For $p\ne q$, the vector spaces $T_pM$ and $T_qM$ have no canonical
identification. A connection supplies a transport rule. The Levi--Civita
connection is characterized by
\[
\nabla g=0,\qquad
\operatorname{Tor}_\nabla=0,
\]
and in coordinates
\[
\Gamma^\rho{}_{\mu\nu}
=\frac12 g^{\rho\sigma}
\left(
\partial_\mu g_{\sigma\nu}
+\partial_\nu g_{\sigma\mu}
-\partial_\sigma g_{\mu\nu}
\right).
\]
The Christoffel symbols are not tensor components. The curvature tensor is
\[
R^\rho{}_{\sigma\mu\nu}
=\partial_\mu\Gamma^\rho{}_{\nu\sigma}
-\partial_\nu\Gamma^\rho{}_{\mu\sigma}
+\Gamma^\rho{}_{\mu\lambda}\Gamma^\lambda{}_{\nu\sigma}
-\Gamma^\rho{}_{\nu\lambda}\Gamma^\lambda{}_{\mu\sigma}.
\]
The Ricci, scalar-curvature, and Einstein tensors are written
\[
\operatorname{Ric}_{\sigma\nu}
=R^\rho{}_{\sigma\rho\nu},
\qquad
\operatorname{Scal}_g
=g^{\sigma\nu}\operatorname{Ric}_{\sigma\nu},
\]
\[
\operatorname{Ein}(g)_{\mu\nu}
=\operatorname{Ric}_{\mu\nu}
-\frac12\operatorname{Scal}_g\,g_{\mu\nu}.
\]
The notation $\operatorname{Ein}(g)$ avoids a collision between the
standard symbol $G_{\mu\nu}$ and the GO-QS symmetry-group symbol $G$.
With geometrized stress-energy, Einstein's equation is
\[
\operatorname{Ein}(g)_{\mu\nu}=8\pi T_{\mu\nu}.
\]
At the SI boundary it is
$\operatorname{Ein}(g)_{\mu\nu}
=(8\pi G_{\rm N}/c^4)T^{\rm SI}_{\mu\nu}$.

An affinely parametrized geodesic satisfies
\[
\nabla_{\dot\gamma}\dot\gamma=0.
\]
Coordinate acceleration is not invariant; the covariant equation is.

\section{Frame/coframe duality and the corrected tetrad formula}

\begin{definition}[Frame and coframe]
A local frame and its dual coframe are
\[
e_a=e_a{}^\mu\partial_\mu,
\qquad
\vartheta^a=e^a{}_\mu\dd x^\mu,
\]
with
\[
e^a{}_\mu e_b{}^\mu=\delta^a{}_b,
\qquad
e^a{}_\mu e_a{}^\nu=\delta_\mu{}^\nu.
\]
The frame is orthonormal when
\[
g_{\mu\nu}e_a{}^\mu e_b{}^\nu=\eta_{ab}.
\]
\end{definition}

\begin{proposition}[Metric reconstruction]
For a Lorentz-orthonormal frame/coframe pair,
\[
\boxed{
g_{\mu\nu}
=\eta_{ab}e^a{}_\mu e^b{}_\nu,
\qquad
g^{\mu\nu}
=\eta^{ab}e_a{}^\mu e_b{}^\nu.
}
\]
\end{proposition}
\begin{proof}
The abstract identity
$g=\eta_{ab}\vartheta^a\otimes\vartheta^b$ gives the first component
formula. Inverting the coframe matrix and the metric gives the second.
Direct contraction with $e_c{}^\mu e_d{}^\nu$ returns $\eta_{cd}$.
\end{proof}

\begin{warning}[Correction to v0.2]
The expression $\eta_{ab}e_a{}^\mu e_b{}^\nu$ reconstructs
$g^{\mu\nu}$, not $g_{\mu\nu}$. The covariant metric requires the dual
coframe coefficients $e^a{}_\mu$.
\end{warning}

\subsection{Example: spherical Minkowski coframe}

On a chart with $r>0$ and $0<\theta<\pi$,
\[
g=-\dd t^2+\dd r^2+r^2\dd\theta^2
+r^2\sin^2\theta\,\dd\phi^2
\]
in $c=1$ units. An orthonormal coframe is
\[
\vartheta^0=\dd t,\quad
\vartheta^1=\dd r,\quad
\vartheta^2=r\,\dd\theta,\quad
\vartheta^3=r\sin\theta\,\dd\phi.
\]
Then $g=\eta_{ab}\vartheta^a\otimes\vartheta^b$. Coordinate singularities
at the polar axis are not curvature singularities.

\section{Observers and transport}

A future-directed observer is a timelike worldline
\[
\gamma:I\to M,\qquad
u=\dot\gamma,\qquad
g(u,u)=-1.
\]
Its acceleration is $a=\nabla_u u$, with $g(a,u)=0$.

\begin{definition}[Observer tetrad]
An observer tetrad along $\gamma$ is an orthonormal frame
$(e_0,e_1,e_2,e_3)$ with $e_0=u$. The spatial axes span
\[
u^\perp=\set{X\in T_{\gamma(\tau)}M:g(X,u)=0}.
\]
\end{definition}

\begin{definition}[Fermi--Walker transport]
In the declared $c=1$ convention, a vector field $X$ along $\gamma$ is
Fermi--Walker transported when
\[
\nabla_uX
=\left(u\otimes a^\flat-a\otimes u^\flat\right)X.
\]
For $a=0$ this reduces to parallel transport.
\end{definition}

In SI variables with $g(u_{\rm SI},u_{\rm SI})=-c^2$, the right-hand side
is divided by $c^2$. This conversion boundary is mandatory.

\begin{proposition}[No canonical spatial triad]
At a point $p$, after fixing a unit future timelike vector $u$, the set of
oriented orthonormal bases of $u^\perp$ is a torsor for $\SO(3)$. Hence
$g$ and $u$ do not select a preferred spatial triad.
\end{proposition}
\begin{proof}
$u^\perp$ is a three-dimensional Euclidean vector space. $\SO(3)$ acts
freely and transitively on its oriented orthonormal bases.
\end{proof}

\subsection{Vacuum is not flatness}

Einstein vacuum, $R_{\mu\nu}=0$, does not imply
$R^\rho{}_{\sigma\mu\nu}=0$. For a Jacobi field $\xi$ along a geodesic,
the declared curvature convention gives
\[
\nabla_u\nabla_u\xi=R(u,\xi)u.
\]
Tidal relative acceleration can therefore provide a local relational
signal even in vacuum.

\section{Frame gauge versus physical apparatus orientation}

The oriented, time-oriented orthonormal frame bundle is
\[
\pi_{\F}:\F^+(M,g)\to M,
\]
with structure group $\SO^+(1,3)$. A local Lorentz gauge transformation
changes a representative frame. A physical apparatus, however, has
material axes, couplings, and a response model. Rotating the apparatus
relative to an external field changes its physical configuration; it is
not erased by relabelling components.

\begin{model}[Gauge-covariant sensor response]
Let $s$ denote physical fields and $F$ a frame representative. A
well-defined scalar readout $\mathcal R$ satisfies
\[
\mathcal R(F\Lambda,\rho(\Lambda)^{-1}s)
=\mathcal R(F,s)
\]
for simultaneous passive frame change. An active apparatus rotation
$A$ acts on the apparatus state alone and may produce
$\mathcal R(A\!\cdot\!F,s)\ne\mathcal R(F,s)$.
\end{model}

\begin{warning}
The phrase ``frame-dependent readout'' is ambiguous unless it states
whether the frame is a gauge representative or a material apparatus
orientation. GO-QS stores these in different protocol fields.
\end{warning}

\section{The GO-QS tensor/GR observation package}

For a relativistic application, use
\[
\mathfrak O_{\rm GR}=
\left(
E,\Phi,\xi,\Lambda,
\set{E_\lambda,Y_\lambda^0,Y_\lambda,
R_\lambda,O_\lambda,K_\lambda,\widehat q_\lambda}_{\lambda\in\Lambda},
G_{\rm diff},G_{\rm LL},\Qreg,\Dprof
\right).
\]
Here $E$ is a state or field configuration space, not spacetime alone.
$G_{\rm diff}$ and $G_{\rm LL}$ record their actions separately. A local
protocol minimally contains
\[
\lambda=
\left(
\gamma,u,\mathcal E_{\rm app},
\varepsilon_x,\delta\tau,B_\nu,
\Sigma_\eta,Q_{\rm quant},
\mathcal C_{\rm coord},\mathcal C_{\rm LL},
T_{\rm hor}
\right).
\]
$\mathcal E_{\rm app}$ is the material apparatus orientation;
$\mathcal C_{\rm LL}$ is a gauge convention.

\section{Pullback geometry and its radical}

Let $O:E\to B$ be smooth and let $h$ be a nondegenerate
pseudo-Riemannian metric on $B$. Define
\[
(O^*h)_e(v,w)
=h_{O(e)}(\dd O_e v,\dd O_e w).
\]
The vertical space is $\mathcal V_e=\ker\dd O_e$.

\begin{theorem}[Exact radical of a pullback form]
\[
\boxed{
\rad((O^*h)_e)
=(\dd O_e)^{-1}
\left(
(\im\dd O_e)^{\perp_h}
\right).
}
\]
In particular, $\ker\dd O_e\subseteq\rad((O^*h)_e)$.
\end{theorem}
\begin{proof}
A vector $v$ belongs to the radical precisely when
$h(\dd O_ev,\dd O_ew)=0$ for every $w\in T_eE$. This is equivalent to
$\dd O_ev$ being orthogonal to $\im\dd O_e$, which gives the stated
preimage.
\end{proof}

\begin{corollary}
If $h$ is Riemannian, then
\[
\rad((O^*h)_e)=\ker\dd O_e.
\]
The same equality holds for pseudo-Riemannian $h$ whenever $O$ is a
submersion at $e$.
\end{corollary}
\begin{proof}
For a positive-definite form,
$\im\dd O_e\cap(\im\dd O_e)^{\perp_h}=\{0\}$. For a submersion,
$\im\dd O_e=T_{O(e)}B$, whose orthogonal complement is zero by
nondegeneracy of $h$.
\end{proof}

\begin{remark}
In Lorentzian signature, injectivity of $\dd O_e$ does not guarantee
nondegeneracy of $O^*h$: the image can be a degenerate null subspace.
Thus ``non-immersion implies degeneracy'' is correct, but its converse
fails.
\end{remark}

\subsection{No automatic Lorentzian distance}

A Lorentzian metric supplies causal cones, proper time on timelike curves,
and a time-separation function under suitable hypotheses. It does not
supply a positive metric distance analogous to a Riemannian distance.
Statistical residuals must use a positive-definite sensor covariance or a
declared loss, not the indefinite form $h$ as a norm.

\section{Descent and recoverability}

\begin{theorem}[Smooth scalar descent]
Let $O:E\to B$ be a surjective submersion and $q:E\to\R$ smooth. There is
a unique smooth $\overline q:B\to\R$ satisfying
\[
q=\overline q\circ O
\]
if and only if $q$ is constant on every fiber $O^{-1}(b)$.
\end{theorem}
\begin{proof}
Necessity is immediate. For sufficiency, define
$\overline q(b)=q(e)$ for $e\in O^{-1}(b)$. A submersion admits smooth
local sections $\sigma:U\to E$, and on $U$,
$\overline q=q\circ\sigma$, proving smoothness. Surjectivity gives
uniqueness.
\end{proof}

\begin{definition}[Horizontal covariant tensor]
A covariant $k$-tensor $T$ on $E$ is horizontal for $O$ if
$T(v_1,\ldots,v_k)=0$ whenever at least one $v_i$ is vertical.
\end{definition}

\begin{theorem}[Descent of covariant tensors]
Let $O:E\to B$ be a surjective submersion with connected fibers, and let
$T$ be a smooth covariant $k$-tensor on $E$. Then there is a unique
covariant tensor $\overline T$ on $B$ with
\[
T=O^*\overline T
\]
if and only if:
\begin{enumerate}[leftmargin=*]
\item $T$ is horizontal; and
\item $\Lie_VT=0$ for every smooth vertical vector field $V$.
\end{enumerate}
\end{theorem}
\begin{proof}[Proof sketch]
Pullbacks are horizontal and invariant under vertical flows. Conversely,
horizontality makes the value of $T$ depend only on projected tangent
vectors. Vertical-flow invariance makes that value constant along each
connected fiber. Local sections of the submersion define compatible local
tensors on $B$, which glue uniquely.
\end{proof}

\begin{warning}[Disconnected fibers and contravariant indices]
If fibers are disconnected, vertical Lie invariance does not compare
different components; explicit cross-component consistency is also
required. The stated criterion is for covariant tensors. Descent of
contravariant or mixed tensors requires a separate projectability
condition and is not obtained by replacing indices formally.
\end{warning}

\section{Readouts on the frame bundle}

\begin{theorem}[Failure to descend to spacetime]
Let $\mathcal R:\F^+(M,g)\to Y$. If two frames $F_1,F_2$ over the same
event satisfy
\[
\pi_{\F}(F_1)=\pi_{\F}(F_2),\qquad
\mathcal R(F_1)\ne\mathcal R(F_2),
\]
then no map $\overline{\mathcal R}:M\to Y$ satisfies
$\mathcal R=\overline{\mathcal R}\circ\pi_{\F}$.
\end{theorem}
\begin{proof}
Such a factorization would assign the same value
$\overline{\mathcal R}(p)$ to both frames, contradicting the hypothesis.
\end{proof}

This set-theoretic theorem does not decide whether the variation is a gauge
artifact or a physical apparatus response. That decision requires the
simultaneous transformation law in the protocol.

\section{Worked examples}

\subsection{Riemannian orthographic projection}

For
\[
\pi:\R^3\to\R^2,\qquad
\pi(x,y,z)=(x,y),
\]
and $h=\dd x^2+\dd y^2$,
\[
\pi^*h=\dd x^2+\dd y^2,\qquad
[\pi^*h]=
\begin{pmatrix}
1&0&0\\
0&1&0\\
0&0&0
\end{pmatrix}.
\]
Its radical equals $\ker\dd\pi=\operatorname{span}\{\partial_z\}$.

\subsection{Lorentzian null immersion}

Let
\[
i:\R\to\R^{1,1},\qquad i(s)=(s,s),
\qquad
h=-\dd t^2+\dd x^2.
\]
$\dd i_s$ is injective, but
\[
i^*h=-\dd s^2+\dd s^2=0.
\]
Thus an immersion can induce a degenerate pullback in indefinite
signature. This is the counterexample excluded by an unjustified
Riemannian intuition.

\section{Quantity and Protocol Ledger}

\subsection{Coordinate-dimension rule}

Store a coordinate passport $[x^\mu]$ for every chart. Since
$[\dd s^2]=L^2$,
\[
[g_{\mu\nu}]
=\frac{L^2}{[x^\mu][x^\nu]}.
\]
If all $x^\mu$ have length dimension, $g_{\mu\nu}$ is dimensionless. If
$x^0=t$ is measured in seconds while spatial coordinates are lengths,
metric components have different dimensions and cannot be treated as an
unannotated numerical matrix.

\begin{longtable}{L{0.20\linewidth}L{0.19\linewidth}L{0.23\linewidth}L{0.279\linewidth}}
\toprule
Object & Abstract dimension & GO-QS kind & Required metadata\\
\midrule
\endhead
$x^\mu$ & coordinate-specific & scalar/angle & chart and unit\\
$g$ & $\dd s^2:L^2$ & covariant tensor & signature and coordinate passport\\
$u$ & unitless in $c=1$ normalization & vector & worldline and parameter\\
$a=\nabla_u u$ & inverse length in geometrized parametrization &
vector & unit context\\
$\Gamma^\rho{}_{\mu\nu}$ & component-specific & connection coefficients &
chart; not a tensor\\
$R^\rho{}_{\sigma\mu\nu}$ & component-specific & tensor &
curvature convention\\
$\mathrm{Scal}_g$ & $L^{-2}$ & scalar & metric and sign convention\\
$e_a{}^\mu$ & coordinate-specific & frame coefficients &
frame, chart, Lorentz gauge\\
$e^a{}_\mu$ & dual coordinate-specific & coframe coefficients &
inverse-pair check\\
$\dd O_e$ & output/input units & linear map & domain, codomain, rank threshold\\
$O^*h$ & inherited from $h$ & covariant tensor & signature and radical\\
\bottomrule
\end{longtable}

\subsection{Minimal relativistic protocol}

\begin{longtable}{L{0.27\linewidth}L{0.64\linewidth}}
\toprule
Field & Required content\\
\midrule
\endhead
\texttt{unit\_context} & SI or GR geometrized; explicit conversion boundary\\
\texttt{signature} & here $(-,+,+,+)$\\
\texttt{curvature\_convention} & definition of $R(X,Y)Z$\\
\texttt{coordinates} & chart, ranges, units, singular sets\\
\texttt{frame} & frame and coframe matrices plus inverse-pair validation\\
\texttt{lorentz\_gauge} & representative and transformation law\\
\texttt{apparatus\_orientation} & material axes, distinct from gauge\\
\texttt{observer} & $\gamma,u,a$ and transport law\\
\texttt{observation} & $R_\lambda,O_\lambda,K_\lambda$ with signatures\\
\texttt{resolution} & spatial, proper-time, spectral, and angular\\
\texttt{uncertainty} & covariance/posterior/risk and calibration\\
\texttt{causal\_domain} & time orientation and admissible $J^-$ region\\
\bottomrule
\end{longtable}

\section{Correction log: v0.2 to v0.3}

\begin{enumerate}[leftmargin=*]
\item Replaced the erroneous covariant metric reconstruction by the
coframe formula and added the inverse-metric frame formula.
\item Declared signature, curvature convention, and
\texttt{GR\_geometrized} unit context.
\item Added the SI normalization $g(u_{\rm SI},u_{\rm SI})=-c^2$ and the
$1/c^2$ Fermi--Walker conversion.
\item Separated passive coordinates, active diffeomorphisms, local Lorentz
gauge, physical apparatus rotation, reduction, and observation.
\item Replaced the informal basic-tensor equivalence by a theorem for a
surjective submersion with connected fibers.
\item Added the disconnected-fiber and mixed-tensor limitations.
\item Computed the exact radical of a pseudo-Riemannian pullback.
\item Added a null-immersion counterexample.
\item Explicitly denied a positive metric-distance interpretation for a
Lorentzian form.
\item Added a coordinate-dimension passport and a complete GO-QS protocol
ledger.
\end{enumerate}

\section{Claim audit}

\textbf{Established statements used:}
\begin{enumerate}[leftmargin=*]
\item tensor contractions are coordinate independent;
\item different tangent spaces require a transport rule;
\item frame and coframe reconstruct inverse and covariant metrics,
respectively;
\item proper acceleration, Fermi--Walker rotation, and curvature are
geometrically meaningful after conventions are fixed;
\item the radical formula characterizes local pullback degeneracy;
\item scalar and covariant-tensor descent require fiber compatibility.
\end{enumerate}

\textbf{Not claimed:}
\begin{enumerate}[leftmargin=*]
\item coordinate transformations are lossy observation channels;
\item every diffeomorphism is gauge under arbitrary boundary conditions;
\item an apparatus orientation is identical to a Lorentz gauge choice;
\item a Lorentzian metric is a positive distance;
\item every hidden tensor is recoverable;
\item this interface modifies general relativity.
\end{enumerate}

\section{Conclusion}

The corrected tensorial interface is summarized by
\[
\boxed{
\begin{gathered}
\text{covariance preserves representation-independent content;}
\\
\text{a declared observation channel may quotient it.}
\end{gathered}
}
\]
The statement is now type safe. Coordinate and gauge changes act
invertibly on representations; apparatus orientation belongs to the
physical protocol; observation maps have explicit domains, codomains, and
kernels; Lorentzian pullbacks are treated as indefinite bilinear forms;
and tensor recovery is a descent problem with stated global hypotheses.

\begin{thebibliography}{99}
\bibitem{KN}
S.~Kobayashi and K.~Nomizu,
\emph{Foundations of Differential Geometry}, Vol. I,
Wiley, 1963.

\bibitem{LeeSmooth}
J.~M. Lee,
\emph{Introduction to Smooth Manifolds}, 2nd ed.,
Springer, 2013.

\bibitem{LeeRiemannian}
J.~M. Lee,
\emph{Introduction to Riemannian Manifolds}, 2nd ed.,
Springer, 2018.

\bibitem{ONeill}
B.~O'Neill,
\emph{Semi-Riemannian Geometry with Applications to Relativity},
Academic Press, 1983.

\bibitem{Wald}
R.~M. Wald,
\emph{General Relativity},
University of Chicago Press, 1984.

\bibitem{Nakahara}
M.~Nakahara,
\emph{Geometry, Topology and Physics}, 2nd ed.,
CRC Press, 2003.

\bibitem{Palais}
R.~S. Palais,
Foundations of global non-linear analysis,
W. A. Benjamin, 1968.
\end{thebibliography}

\end{document}
